\section{A Scalarized Mixing Bound}
\label{app:beta}
The exact constrained frontier in Theorem~\ref{thm:limit} is binary and Markov.
This appendix gives a different statement for a general stationary capacity
process. It bounds the value of delayed information for one bounded scalar
objective. It does not establish the same matched-mismatch frontier for arbitrary
capacity.

Let $\{a_t\}$ be stationary on $\{0,\dots,C\}$, with
$\mathcal F_{\le s}=\sigma(a_u:u\le s)$ and
$\mathcal F_{\ge s}=\sigma(a_u:u\ge s)$. Define
\begin{equation}\label{eq:beta}
\beta(D)=\mathbb E\!\left[
\sup_{B\in\mathcal F_{\ge t}}
\left|\Pr(B\mid\mathcal F_{\le t-D})-\Pr(B)\right|
\right].
\end{equation}
Fix a bounded per-epoch objective $R(w,a)$ with oscillation
$\Omega=\max_{w,a}R(w,a)-\min_{w,a}R(w,a)$. One example is the scalarized
objective $R_\lambda(w,a)=\min(w,a)-\lambda(w-a)^+$. Let $V_R^\star(D)$ be the
best value of a lag-$D$ causal policy and let $V_R^{\mathrm{stat}}$ be the best
value of a record-independent action.

\begin{theorem}[$\beta$-mixing value bound]\label{thm:beta}
For any such process and objective,
\begin{equation}\label{eq:betabound}
0\le V_R^\star(D)-V_R^{\mathrm{stat}}\le\Omega\beta(D).
\end{equation}
Thus the scalarized value of delayed information vanishes for an absolutely
regular process and decays geometrically when $\beta(D)$ does.
\end{theorem}
\begin{proof}
Let $\mathcal G=\mathcal F_{\le t-D}$,
$\mu_{\mathcal G}=\Pr(a_t\in\cdot\mid\mathcal G)$, and
$\psi(\mu)=\max_w\sum_a R(w,a)\mu(a)$. Pointwise optimization gives
$V_R^\star(D)=\mathbb E[\psi(\mu_{\mathcal G})]$ and
$V_R^{\mathrm{stat}}=\psi(\bar\mu)$. Convexity of $\psi$ gives the lower bound.
The function $\psi$ is $\Omega$-Lipschitz in total variation because it is the
maximum of linear functionals whose ranges are at most $\Omega$. Hence
\[
V_R^\star(D)-V_R^{\mathrm{stat}}
\le \Omega\,\mathbb E\|\mu_{\mathcal G}-\bar\mu\|_{\mathrm{TV}}
\le \Omega\beta(D),
\]
where the last inequality follows because $\sigma(a_t)$ is contained in
$\mathcal F_{\ge t}$.
\end{proof}

For the symmetric two-state Markov process,
$\beta(D)=|1-2q(D)|/2=r(D)/2$ under the convention in \eqref{eq:beta}. This shows
that the exponential decay rate in Theorem~\ref{thm:limit} is a tight mixing-rate
example. Equality of the scalarized bound with the constrained vertical frontier
requires the binary objective geometry proved in Section~\ref{sec:limit}.

\section{Details of the Branch-Selection Rate}
\label{app:regret}
Algorithm~\ref{alg:adaptive} observes delayed but clean samples
$y_i=x_{i-D_{\mathrm{obs}}}$. The fixed shift does not change adjacent flips, so
\[
\hat p_n=\frac1n\sum_{i=1}^n\mathbf 1[y_i\neq y_{i-1}]
\]
is the average of independent Bernoulli$(p)$ variables under the symmetric
Markov model. The controller projects $\hat p_n$ onto $[0,1/2]$ before computing
$\hat r_n=(1-2\hat p_n)^D$.

\noindent\textbf{Upper bound.}
In a local neighborhood of $p_\theta$, binomial Bernstein concentration gives
$|\hat p_n-p|=O(\sigma_\theta/\sqrt n)$ at constant confidence. The derivative
of $r(p)=(1-2p)^D$ has magnitude $s_\theta$ at $p_\theta$. The branch-value
difference satisfies
$d(r)=d'(\theta)(r-\theta)+o(|r-\theta|)$, so a wrong branch incurs local loss of
order $L|r-\theta|$. Summing the resulting $n^{-1/2}$ bound over
$n=1,\dots,N$ and dividing by $N$ gives
$O(Ls_\theta\sigma_\theta/\sqrt N)$.

\noindent\textbf{Lower bound.}
Take two flip probabilities
$p_{\pm}=p_\theta\pm\delta/2$. Their $N$-sample Kullback-Leibler divergence is
$N\delta^2/(2\sigma_\theta^2)+O(N\delta^3)$. Choosing
$\delta=c\sigma_\theta/\sqrt N$ for a sufficiently small constant $c$ keeps the
two laws statistically indistinguishable with constant probability. Their
correlations differ by
$s_\theta\delta+o(\delta)$ and their better branches are opposite. Le Cam's
two-point lemma \cite{tsybakov} therefore gives
$\Omega(Ls_\theta\sigma_\theta/\sqrt N)$ average loss on one of the two
processes.

\noindent\textbf{Exact sensitivity and asymptotics.}
Let $u=\theta^{1/D}=1-2p_\theta$. Since
$\sigma_\theta=\tfrac12\sqrt{1-u^2}$ and
$s_\theta=2D u^{D-1}$,
\begin{equation}\label{eq:exactsens}
s_\theta\sigma_\theta
=D u^{D-1}\sqrt{1-u^2}
=D\theta^{(D-1)/D}\sqrt{1-\theta^{2/D}}.
\end{equation}
For fixed $\theta\in(0,1)$ and $D\to\infty$,
$u=1-\log(1/\theta)/D+o(D^{-1})$, which gives
\[
s_\theta\sigma_\theta
\sim \theta\sqrt{2D\log(1/\theta)}.
\]
For $D\ge2$, differentiating
$(D-1)\log u+\tfrac12\log(1-u^2)$ gives the maximizer
$u^2=(D-1)/D$. Substitution yields \eqref{eq:sup} and the switch level
$\theta=((D-1)/D)^{D/2}$.

\noindent\textbf{Fixed alternatives.}
When $p$ is separated from $p_\theta$ by a constant, Chernoff concentration makes
the probability of choosing the wrong branch decay exponentially in the
cumulative sample size. The expected total branch loss is then finite, so the
average loss over horizon $N$ is $O(1/N)$. A fixed sliding window does not have
this property because its error probability does not vanish.
